%--------------------------------------------------------------------------
%--------------------- Resumo em Português --------------------------------
%--------------------------------------------------------------------------

%\setlength{\absparsep}{18pt} % ajusta o espaçamento dos parágrafos do resumo
\begin{resumo}
Neste relatório técnico serão detalhadas as técnicas utilizadas para a obtenção de dados de plataforma ITSM, bem como as escolhas necessárias para a montagem do painel de visualização de dados. O trabalho foi realizado no âmbito do estágio de Análise de Sistemas na EDF Brasil e as atividades foram feitas em regime presencial. As responsabilidades durante o período de estágio concentraram-se em manutenção de sistemas utilizados pela empresa, fornecimento de dados através das soluções vigentes e montagem e análise de relatórios para as áreas clientes.
\textcolor{red}{(Ficou resumido demais. O ideal é o dobro desse tamanho. Mas é melhor deixarmos para pensar nele mais pro final da escrita.)}

 \vspace{\onelineskip}
 \noindent
 \textbf{Palavras-chave}: Análise de Dados; REST; \textit{Power BI}, Power Query, DAX.
\end{resumo}

%--------------------------------------------------------------------------
%--------------------- Resumo em Inglês --------------------------------
%--------------------------------------------------------------------------
\begin{resumo}[ABSTRACT]
\begin{otherlanguage*}{english}
The abstract text must go here and have the paragraph tab, on the first line,
single-spaced and must contain 150 to 500 words. It should not exceed one sheet.
Attention: it does not have paragraph indentation.In the current world, the execution of program points brings about a process of reformulation and modernization of the required financial and administrative conditions. We can already envision how the consensus on the need for qualification indirectly impacts the reassessment of the sectoral budget. In this way, the perception of difficulties is part of a process of managing the intended indexes.
The commitment to analyzing the constant dissemination of information challenges the ability to balance alternatives to orthodox solutions. Clearly, the commitment among teams has not yet convincingly demonstrated that it will participate in changing corporate paradigms. Nonetheless, the current structure of the organization prepares us to face atypical situations resulting from forms of action. Therefore, the challenging globalized scenario offers an interesting opportunity to verify development guidelines for the future. On the other hand, the new structural model advocated here presents trends towards approving the maintenance of normative conduct rules.

\vspace{\onelineskip}
\noindent 
\textbf{Keywords}: sleep apnea syndrome; placebo; polysomnography; \textit{Aedes aegypti}.
\end{otherlanguage*}
\end{resumo}